\clearpage
\section*{Annexe --- Les trente-cinq trades}
\addcontentsline{toc}{section}{Annexe --- Les trente-cinq trades}

Extrait du journal d'exécution du backtest publié. Les prix sont ceux
réellement obtenus par le courtier~; le score et le levier proviennent du
journal du simulateur au moment de la décision. Le rendement de prix est
calculé de prix d'entrée à prix de sortie, hors portage. Le résultat net inclut
le portage et la commission.

\medskip

{\footnotesize
% Généré par scripts/build_gold_trades_figures.py — NE PAS ÉDITER À LA MAIN.
\begin{longtable}{@{}rllrrrrrrrr@{}}
\caption{Les 35 trades de la sleeve or, backtest MT5 de production}\\
\toprule
\# & Entrée & Sortie & j & Lots & Lev. & Prix in & Prix out & \% prix & Swap & Net \\
\midrule
\endfirsthead
\toprule
\# & Entrée & Sortie & j & Lots & Lev. & Prix in & Prix out & \% prix & Swap & Net \\
\midrule
\endhead
1 & 2021-11-09 & 2021-11-24 & 15 & 0.16 & 5.0 & 1\,832.5 & 1\,788.1 & -2.4 & -141 & -852 \\
2 & 2021-12-08 & 2021-12-09 & 1 & 0.16 & 5.3 & 1\,785.7 & 1\,775.3 & -0.6 & -28 & -195 \\
3 & 2021-12-17 & 2021-12-22 & 5 & 0.19 & 6.5 & 1\,796.3 & 1\,804.8 & +0.5 & -34 & +128 \\
4 & 2021-12-27 & 2022-01-07 & 11 & 0.19 & 6.6 & 1\,812.8 & 1\,795.8 & -0.9 & -146 & -470 \\
5 & 2022-01-10 & 2022-01-17 & 7 & 0.17 & 6.1 & 1\,801.1 & 1\,818.7 & +1.0 & -70 & +230 \\
6 & 2022-01-18 & 2022-05-02 & 104 & 0.18 & 6.4 & 1\,813.9 & 1\,862.7 & +2.7 & -1\,103 & -226 \\
7 & 2022-11-24 & 2023-02-21 & 89 & 0.11 & 3.1 & 1\,755.4 & 1\,834.6 & +4.5 & -564 & +308 \\
8 & 2023-02-22 & 2023-02-23 & 1 & 0.17 & 4.2 & 1\,824.9 & 1\,823.3 & -0.1 & -30 & -57 \\
9 & 2023-03-03 & 2023-03-08 & 5 & 0.26 & 6.6 & 1\,855.3 & 1\,813.3 & -2.3 & -46 & -1\,138 \\
10 & 2023-03-13 & 2023-05-29 & 77 & 0.16 & 4.8 & 1\,913.7 & 1\,942.5 & +1.5 & -726 & -267 \\
11 & 2023-06-01 & 2023-06-05 & 4 & 0.18 & 5.4 & 1\,977.8 & 1\,961.6 & -0.8 & -21 & -312 \\
12 & 2023-06-07 & 2023-06-08 & 1 & 0.16 & 5.0 & 1\,940.2 & 1\,965.4 & +1.3 & -28 & +376 \\
13 & 2023-06-09 & 2023-06-12 & 3 & 0.15 & 4.5 & 1\,960.9 & 1\,957.3 & -0.2 & -9 & -62 \\
14 & 2023-06-13 & 2023-06-14 & 1 & 0.18 & 5.2 & 1\,943.8 & 1\,941.7 & -0.1 & -11 & -48 \\
15 & 2023-07-13 & 2023-07-17 & 4 & 0.22 & 6.6 & 1\,960.6 & 1\,954.7 & -0.3 & -26 & -156 \\
16 & 2023-07-19 & 2023-07-24 & 5 & 0.22 & 6.6 & 1\,977.3 & 1\,954.6 & -1.1 & -65 & -563 \\
17 & 2023-07-25 & 2023-07-28 & 3 & 0.22 & 6.6 & 1\,965.2 & 1\,959.0 & -0.3 & -65 & -202 \\
18 & 2023-07-31 & 2023-08-02 & 2 & 0.20 & 6.2 & 1\,965.7 & 1\,934.7 & -1.6 & -24 & -644 \\
19 & 2023-08-07 & 2023-08-08 & 1 & 0.20 & 6.6 & 1\,936.8 & 1\,925.0 & -0.6 & -12 & -248 \\
20 & 2023-10-19 & 2024-01-25 & 98 & 0.10 & 3.5 & 1\,974.5 & 2\,020.0 & +2.3 & -577 & -122 \\
21 & 2024-01-26 & 2024-02-09 & 14 & 0.21 & 6.6 & 2\,018.7 & 2\,024.9 & +0.3 & -173 & -43 \\
22 & 2024-02-20 & 2024-02-22 & 2 & 0.19 & 6.6 & 2\,024.5 & 2\,024.6 & +0.0 & -45 & -43 \\
23 & 2024-02-26 & 2024-02-28 & 2 & 0.19 & 6.6 & 2\,032.5 & 2\,033.7 & +0.1 & -22 & +0 \\
24 & 2024-02-29 & 2024-06-18 & 110 & 0.19 & 6.6 & 2\,044.3 & 2\,329.3 & +13.9 & -1\,210 & +4\,206 \\
25 & 2024-06-19 & 2024-06-26 & 7 & 0.11 & 3.3 & 2\,328.5 & 2\,298.0 & -1.3 & -45 & -381 \\
26 & 2024-07-03 & 2024-12-04 & 154 & 0.16 & 4.8 & 2\,356.3 & 2\,650.0 & +12.5 & -1\,452 & +3\,246 \\
27 & 2024-12-10 & 2024-12-19 & 9 & 0.11 & 3.1 & 2\,692.0 & 2\,596.8 & -3.5 & -71 & -1\,119 \\
28 & 2024-12-30 & 2025-01-07 & 8 & 0.14 & 4.1 & 2\,607.9 & 2\,650.4 & +1.6 & -66 & +529 \\
29 & 2025-01-13 & 2025-07-31 & 199 & 0.24 & 6.6 & 2\,662.4 & 3\,290.1 & +23.6 & -2\,844 & +12\,219 \\
30 & 2025-08-04 & 2025-08-20 & 16 & 0.21 & 4.0 & 3\,373.8 & 3\,348.4 & -0.8 & -198 & -732 \\
31 & 2025-08-21 & 2026-03-23 & 214 & 0.23 & 4.4 & 3\,338.8 & 4\,405.1 & +31.9 & -2\,874 & +21\,652 \\
32 & 2026-04-01 & 2026-04-02\,$\dagger$ & 0 & 0.10 & 1.7 & 4\,758.8 & 4\,565.1 & -4.1 & -18 & -1\,955 \\
33 & 2026-04-02 & 2026-04-06 & 4 & 0.10 & 1.7 & 4\,676.4 & 4\,649.7 & -0.6 & -12 & -280 \\
34 & 2026-04-15 & 2026-04-17 & 2 & 0.09 & 1.7 & 4\,790.9 & 4\,834.5 & +0.9 & -21 & +371 \\
35 & 2026-04-20 & 2026-04-21 & 1 & 0.10 & 1.8 & 4\,821.0 & 4\,720.7 & -2.1 & -6 & -1\,008 \\
\bottomrule
\multicolumn{11}{@{}l@{}}{\footnotesize $\dagger$ sortie déclenchée par le stop de sécurité, hors borne de séance.}
\end{longtable}

}

\bigskip

\subsection*{Reproduire ces chiffres}

\begin{enumerate}[leftmargin=1.4em, itemsep=0.3em]
\item \code{python scripts/extract\_gold\_trades.py --expect 35} --- relit le
      journal des transactions du simulateur et celui du testeur, apparie les
      positions, puis écrit le tableau des trades et sa fiche de provenance
      dans \code{reports/mt5/}.
\item \code{python scripts/build\_gold\_trades\_figures.py} --- produit les
      figures et les tables de ce document.
\item \code{scripts/compile\_latex\_report.sh goldtrades} --- compile le
      présent document.
\end{enumerate}

La fiche de provenance \code{reports/mt5/gold\_trades\_production.json}
enregistre l'empreinte du fichier source et vérifie qu'il s'agit bien du
backtest publié dans le rapport d'investissement. L'extraction échoue plutôt
que de produire un tableau incomplet si le nombre de positions ne correspond
pas ou si le journal du testeur ne contient pas le run attendu.
